% Nội dung tiếng Việt; tiêu đề giữ tiếng Anh.
% Số liệu tổng hợp từ 48 phản hồi khảo sát do nhóm cung cấp.
\unnumberedchapter{INTRODUCTION}
\label{chap:introduction}

\section{MOTIVATION}

Một người học có thể nhớ nghĩa của một từ tiếng Nhật nhưng vẫn lúng túng khi dùng từ đó để trả lời một câu hỏi. Họ có thể hiểu mẫu ngữ pháp trong sách, nhưng khi nghe một đoạn hội thoại lại không nhận ra mẫu câu vừa học. Khoảng cách giữa việc ghi nhớ kiến thức và sử dụng được ngôn ngữ là vấn đề mà đề tài NihonGO hướng đến. Để thu hẹp khoảng cách này, người học cần có cơ hội luyện tập, được giải thích những phần chưa hiểu và thường xuyên vận dụng lại kiến thức đã học.

Trong bối cảnh chuyển đổi số trong giáo dục, việc tự học có thể diễn ra qua ứng dụng di động, video, từ điển trực tuyến và các công cụ AI. Những hình thức này giúp người học chủ động lựa chọn thời gian và nội dung phù hợp với bản thân. Tuy nhiên, có nhiều tài liệu và công cụ để lựa chọn cũng đặt ra một câu hỏi thực tế: người học nên bắt đầu từ đâu, ôn lại như thế nào và làm sao biết mình đang tiến bộ? Với sinh viên phải sắp xếp việc học ngoại ngữ cùng các môn chuyên ngành, sự thuận tiện và rõ ràng trong quá trình học có ý nghĩa thiết thực.

Tiếng Nhật còn đặt ra yêu cầu kết hợp nhiều loại kiến thức. Bên cạnh Hiragana và Katakana, người học cần ghi nhớ Kanji, cách đọc, từ vựng và các mẫu ngữ pháp. Những kiến thức này tiếp tục xuất hiện trong cả bốn kỹ năng nghe, nói, đọc và viết. Vì vậy, một hệ thống hỗ trợ học tiếng Nhật cần giúp người học thấy được mối liên hệ giữa điều vừa học và cách sử dụng trong ngữ cảnh cụ thể. Đây là hướng tiếp cận được nhóm lựa chọn khi xây dựng NihonGO.

\subsection{Learner Needs from the Survey}

Để tìm hiểu những khó khăn thực tế, nhóm sử dụng dữ liệu khảo sát gồm 48 phản hồi được ghi nhận từ ngày 16/09/2026 đến ngày 17/11/2026. Nội dung khảo sát xoay quanh trình độ hiện tại, khó khăn khi học tiếng Nhật, công cụ đang sử dụng và những chức năng người học mong muốn. Trong số người tham gia, 47 người là sinh viên và một người học vì sở thích. Về trình độ tự đánh giá, 37 người ở mức N5--N4, hai người vừa bắt đầu, sáu người ở mức N3 và ba người ở mức N2--N1. Nhóm mới bắt đầu và N5--N4 chiếm 81,3\%, cho thấy nhu cầu của người học ở giai đoạn xây dựng nền tảng là trọng tâm đáng quan tâm trong mẫu khảo sát.

\begin{table}[htbp]
  \centering
  \caption{Các khó khăn đối với người bắt đầu học tiếng Nhật được lựa chọn trong khảo sát}
  \label{tab:initial-learning-difficulties}
  \begin{tabular}{@{}p{10cm}rr@{}}
    \toprule
    \textbf{Khó khăn} & \textbf{Số người} & \textbf{Tỷ lệ} \\
    \midrule
    Ghi nhớ nét viết và ý nghĩa Kanji & 36 & 75,0\% \\
    Duy trì động lực và tần suất học tập & 35 & 72,9\% \\
    Xây dựng vốn từ vựng nền tảng & 20 & 41,7\% \\
    Phát âm chính xác và trọng âm & 13 & 27,1\% \\
    Học và ghi nhớ Hiragana, Katakana & 10 & 20,8\% \\
    \bottomrule
  \end{tabular}
  \par\smallskip
  {\fontsize{11pt}{14pt}\selectfont Nguồn: tổng hợp từ 48 phản hồi khảo sát của đề tài. Câu hỏi cho phép chọn nhiều đáp án; tỷ lệ được tính trên 48 người trả lời.\par}
\end{table}

Bảng~\ref{tab:initial-learning-difficulties} trình bày ý kiến của toàn bộ người tham gia về khó khăn đối với người bắt đầu học. Kanji và việc duy trì động lực là hai nội dung được lựa chọn nhiều nhất. Điều này gợi ý rằng sản phẩm cần quan tâm đồng thời đến cách học và trải nghiệm học. Người học cần được hỗ trợ ghi nhớ kiến thức, nhưng cũng cần có lý do để quay lại luyện tập mỗi ngày và nhận thấy kết quả từ những lần luyện tập đó.

Các câu trả lời mở giúp làm rõ hơn những con số trên. Có người mong muốn học Kanji qua bộ thủ và câu chuyện để dễ liên tưởng; có người muốn lưu kiến thức đã học vào một nơi thuận tiện để tìm lại. Một số ý kiến đề nghị bổ sung ví dụ từ video, bài hát hoặc những câu nói quen thuộc để việc học từ vựng bớt đơn điệu. Những mong muốn này cho thấy người học quan tâm đến cách kiến thức được giải thích, ghi nhớ và ôn lại, bên cạnh số lượng bài học được cung cấp.

Một vấn đề khác là thiếu cơ hội giao tiếp. Trong khảo sát, có người cho biết đã học từ vựng và ngữ pháp nhưng khi nói vẫn chưa tạo được một câu hoàn chỉnh. Những người khác muốn trò chuyện với AI, nghe lại bản thu hoặc được góp ý về cách dùng từ. Có phản hồi nêu khó khăn khi tự luyện phát âm qua video: người học nghe và bắt chước nhưng không biết mình sai ở đâu. Các ý kiến này là cơ sở để nhóm chú trọng hoạt động thực hành và phản hồi trong thiết kế sản phẩm.

Khảo sát cũng ghi nhận việc sử dụng nhiều công cụ: 36 người chọn Mazii, 30 người chọn Quizlet, 19 người chọn Anki và 45 người chọn ChatGPT, Gemini hoặc công cụ trò chuyện AI khác. Người trả lời được chọn nhiều phương án, nên một người có thể sử dụng đồng thời nhiều công cụ. Một số phản hồi đánh giá các ứng dụng hiện có hữu ích cho tra cứu và ghi nhớ; một số khác vẫn mong muốn được hỗ trợ thêm về giao tiếp, lộ trình và ôn tập. Các ý kiến về chi phí cũng gợi mở yêu cầu cân nhắc khả năng tiếp cận của sinh viên khi đề xuất mô hình kinh doanh.

Với quy mô 48 phản hồi và thành phần chủ yếu là sinh viên, khảo sát được dùng để định hướng sản phẩm cho nhóm người dùng ban đầu. Trình độ trong khảo sát là thông tin người học tự đánh giá. Những kết quả này chưa đại diện cho toàn bộ thị trường học tiếng Nhật, nhưng cung cấp các vấn đề cụ thể để nhóm lựa chọn và giải quyết trong phạm vi đồ án.

\subsection{The Direction of NihonGO}

Từ những nhu cầu trên, NihonGO được định hướng trở thành một môi trường học tập kết hợp kiến thức nền tảng với thực hành bốn kỹ năng. Chẳng hạn, khi gặp một từ mới trong bài đọc, người học cần hiểu nghĩa và cách dùng của từ; sau đó có thể ôn lại bằng flashcard và thử vận dụng trong bài viết hoặc hội thoại. Sự kết nối này là mục tiêu thiết kế của hệ thống, giúp mỗi hoạt động học có ý nghĩa rõ ràng trong quá trình sử dụng tiếng Nhật.

AI được lựa chọn để hỗ trợ những hoạt động cần tương tác, như trò chuyện theo tình huống, giải thích một phần nội dung bài đọc hoặc góp ý bài viết. Các thành phần nhận dạng và tổng hợp tiếng nói hỗ trợ chuyển đổi giữa âm thanh và văn bản, tạo điều kiện xây dựng hoạt động luyện nói và nghe lại. Ví dụ, với một tình huống gọi món, người học có thể thử lựa chọn cách diễn đạt, tiếp nhận câu trả lời và tiếp tục cuộc hội thoại. Cách tổ chức này hướng đến việc tạo thêm cơ hội luyện tập ngoài lớp học.

Bên cạnh các chức năng AI, hệ thống cần có học liệu, công cụ tra cứu và cách lưu kết quả rõ ràng. Đây là những thành phần giúp người học tiếp tục ôn tập sau một lần tương tác. Từ định hướng đó, nhóm lựa chọn đề tài \textit{Developing an AI-Integrated System to Support the Development of Four Japanese Language Skills -- NihonGO}. Đề tài tập trung xây dựng một sản phẩm có thể dùng thử, đánh giá và tiếp tục hoàn thiện từ trải nghiệm thực tế của người học.

\section{OBJECTIVES AND CONTRIBUTIONS}

\subsection{General Objective}

Mục tiêu tổng quát của đề tài là xây dựng NihonGO, hệ thống tích hợp AI hỗ trợ phát triển bốn kỹ năng nghe, nói, đọc và viết tiếng Nhật. Hệ thống hướng đến người học tiếng Nhật ở mọi trình độ, có nhu cầu tự học hoặc luyện tập bổ sung bên cạnh các khóa học. Thông qua việc kết hợp bài luyện, tra cứu, ôn tập và phản hồi, NihonGO tạo điều kiện để người học củng cố kiến thức và thử vận dụng vào những tình huống cụ thể.

Trong phạm vi đồ án, nhóm hướng đến một phiên bản sản phẩm khả dụng tối thiểu (Minimum Viable Product -- MVP). Đây là phiên bản có các chức năng cốt lõi đủ để người học trải nghiệm quá trình học và để nhóm kiểm tra mức độ phù hợp của giải pháp. Kết quả dùng thử sẽ làm cơ sở hoàn thiện sản phẩm, lựa chọn chức năng ưu tiên và đề xuất hướng phát triển thương mại sau đồ án.

\subsection{Specific Objectives}

Các mục tiêu cụ thể được xác định theo quá trình xây dựng sản phẩm:

\begin{enumerate}[leftmargin=*, itemsep=4pt]
  \item \textbf{Phân tích nhu cầu người học.} Tổng hợp dữ liệu khảo sát, xác định những khó khăn nổi bật và chuyển thành yêu cầu của hệ thống. Kết quả cần làm rõ sản phẩm phục vụ ai, hỗ trợ hoạt động nào và chức năng nào cần được ưu tiên.
  \item \textbf{Thiết kế hệ thống.} Xây dựng kiến trúc, cơ sở dữ liệu và giao diện cho các hoạt động học tập. Thiết kế cần giúp người dùng dễ tìm bài luyện, thực hiện nhiệm vụ và xem lại kết quả; đồng thời thuận tiện cho việc bảo trì, bổ sung học liệu và mở rộng chức năng.
  \item \textbf{Xây dựng phiên bản MVP.} Phát triển các chức năng luyện nghe, hội thoại theo tình huống, đọc hiểu và luyện viết; kết hợp tra cứu từ vựng, Kanji, flashcard và theo dõi học tập. Mỗi chức năng cần có một quá trình sử dụng rõ ràng từ khi bắt đầu đến khi nhận kết quả.
  \item \textbf{Tích hợp AI vào hoạt động học.} Sử dụng AI để hỗ trợ hội thoại, giải thích và góp ý theo nội dung người học đang luyện tập. Phản hồi cần dễ hiểu, liên quan đến bài học và được trình bày sao cho người học biết phần nào cần xem lại hoặc tiếp tục thực hành.
  \item \textbf{Kiểm thử và đánh giá sản phẩm.} Kiểm tra độ ổn định của các chức năng, khả năng lưu dữ liệu, thời gian phản hồi và cách xử lý lỗi. Kết hợp dùng thử để tìm hiểu người học có hoàn thành được hoạt động chính hay không và còn gặp trở ngại ở bước nào.
  \item \textbf{Đề xuất mô hình kinh doanh.} Xác định giá trị của NihonGO đối với người dùng, nhóm khách hàng ban đầu và phương án duy trì sản phẩm. Việc đề xuất gói sử dụng cần cân nhắc cả chi phí vận hành AI lẫn khả năng chi trả của sinh viên.
\end{enumerate}

\subsection{Expected Contributions}

Kết quả dự kiến của đề tài gồm bản phân tích nhu cầu, thiết kế hệ thống và phiên bản NihonGO có thể triển khai dùng thử. Điểm đóng góp của sản phẩm nằm ở cách tổ chức các hoạt động học tiếng Nhật có liên hệ với nhau: từ tiếp cận kiến thức, luyện tập đến xem lại và ôn tập. Các nhu cầu được nêu trong khảo sát, như khó ghi nhớ Kanji, thiếu cơ hội nói và khó duy trì việc học, là cơ sở để giải thích các lựa chọn chức năng.

Về kỹ thuật, đề tài vận dụng các công nghệ sẵn có để xây dựng một hệ thống học tập hoàn chỉnh ở quy mô MVP. Nhóm tập trung vào cách kết hợp ứng dụng, dữ liệu học tập và dịch vụ AI để phục vụ nhiệm vụ cụ thể của người học. Quá trình kiểm thử và dùng thử giúp ghi nhận những điểm làm được, những hạn chế còn tồn tại và các phần cần ưu tiên khi phát triển phiên bản tiếp theo.

\section{SCOPE OF THE THESIS}

\subsection{Target Users}

NihonGO hướng đến người học tiếng Nhật ở tất cả các trình độ, từ người mới bắt đầu đến người học ở các mức N5, N4, N3, N2 và N1. Đối tượng sử dụng bao gồm sinh viên, người đi làm và người tự học vì sở thích hoặc mục tiêu cá nhân. Tùy theo trình độ và nhu cầu, người học có thể lựa chọn nội dung để xây dựng kiến thức nền tảng, củng cố kỹ năng hoặc luyện cách sử dụng tiếng Nhật trong những ngữ cảnh phức tạp hơn.

Hệ thống hỗ trợ cả việc tự học độc lập và luyện tập bổ sung bên cạnh các khóa học. Người mới bắt đầu có thể làm quen với chữ viết, từ vựng và mẫu câu cơ bản; người đã có nền tảng có thể luyện nghe, hội thoại, đọc hiểu và diễn đạt bằng văn bản theo trình độ của mình. Dù mẫu khảo sát chủ yếu là sinh viên ở mức N5--N4, đối tượng phục vụ của NihonGO được xác định rộng hơn, nhằm đáp ứng nhu cầu học tập xuyên suốt từ cơ bản đến nâng cao.

\subsection{Functional and Technology Scope}

Phạm vi chức năng của hệ thống được tổ chức theo bốn kỹ năng, với các công cụ kiến thức nền tảng dùng chung:

\begin{itemize}[leftmargin=*, itemsep=4pt]
  \item \textbf{Nghe:} luyện nghe với nội dung âm thanh hoặc video có văn bản hỗ trợ; cho phép theo dõi nội dung và thực hành nghe, nhắc lại theo từng phần.
  \item \textbf{Nói:} luyện hội thoại với AI theo tình huống, sử dụng kiến thức đã học để tạo câu trả lời và xem lại nội dung tương tác.
  \item \textbf{Đọc:} đọc văn bản, trả lời câu hỏi và xem giải thích gắn với nội dung bài; hỗ trợ tìm hiểu từ vựng và cấu trúc câu khi gặp phần chưa hiểu.
  \item \textbf{Viết:} luyện diễn đạt bằng văn bản, lưu bản nháp hoặc bài nộp và tiếp nhận góp ý để xem xét chỉnh sửa.
  \item \textbf{Tra cứu và ôn tập:} cung cấp công cụ tra từ, tìm hiểu Kanji và học bằng flashcard; lưu thông tin học tập để người dùng thuận tiện tiếp tục ở những lần sau.
\end{itemize}

Về công nghệ, ứng dụng sử dụng React Native và TypeScript; dịch vụ máy chủ được xây dựng bằng Flask và dữ liệu học tập được quản lý bằng PostgreSQL. Phiên bản web cho phép truy cập bằng trình duyệt. AI và các thành phần xử lý tiếng nói được tích hợp theo yêu cầu của từng chức năng. Cách lựa chọn và kết nối các thành phần này được trình bày chi tiết ở Chương 2.

\subsection{System Limitations}

Trong phạm vi đồ án, nhóm ưu tiên hoàn thiện các hoạt động học cốt lõi trên bộ học liệu được chuẩn bị. Nội dung và mức độ hỗ trợ có thể khác nhau giữa các trình độ. Những mong muốn mở rộng từ khảo sát, như kết nối trực tiếp với người bản xứ, gọi video hoặc xây dựng lộ trình tự động cho từng cá nhân, được xem xét cho các giai đoạn phát triển sau.

Phản hồi của AI được dùng để hỗ trợ người học và có thể có sai sót. Đặc biệt, việc chuyển giọng nói thành văn bản chưa đủ để kết luận người học phát âm đúng hoặc sai ở từng âm, nên khả năng chấm phát âm cần được đánh giá riêng. Các hoạt động phụ thuộc AI và dữ liệu trên máy chủ cần kết nối Internet. Hệ thống không thay thế vai trò hướng dẫn của giáo viên và không cung cấp chứng nhận năng lực tiếng Nhật.

Đánh giá trong đồ án tập trung vào chức năng, độ ổn định và trải nghiệm dùng thử. Việc xác định người học tiến bộ bao nhiêu sau một thời gian sử dụng cần có quá trình theo dõi và kiểm tra đầu vào, đầu ra phù hợp. Vì vậy, kết quả khảo sát nhu cầu ban đầu được dùng để định hướng thiết kế, không được xem là bằng chứng về hiệu quả học tập của sản phẩm.

\subsection{Deployment Environment}

Hệ thống được phát triển và kiểm thử trên máy cục bộ trước khi đưa lên môi trường dùng thử qua Internet. Người dùng truy cập bằng ứng dụng hoặc trình duyệt; các hoạt động hội thoại bằng giọng nói cần thiết bị có micro và quyền ghi âm. Máy chủ tiếp nhận yêu cầu, xử lý dữ liệu học tập và kết nối đến các dịch vụ AI cần thiết.

Môi trường triển khai cần hỗ trợ quản lý tài khoản, bảo vệ dữ liệu và lưu kết quả theo đúng người dùng. Khóa truy cập dịch vụ AI được quản lý ở phía máy chủ. Trong quá trình dùng thử, nhóm sẽ theo dõi lỗi, thời gian phản hồi và chi phí để điều chỉnh cấu hình phù hợp với quy mô sử dụng. Đây cũng là cơ sở để xem xét khả năng mở rộng khi sản phẩm có thêm người học.

\section{RESEARCH METHODS}

Đề tài được thực hiện theo quá trình khảo sát, thiết kế, xây dựng và đánh giá. Trước hết, nhóm tổng hợp câu trả lời để xác định người dùng đang gặp khó khăn gì và mong muốn được hỗ trợ như thế nào. Các câu hỏi lựa chọn được thống kê theo số người và tỷ lệ; câu trả lời mở được đọc và nhóm theo những nội dung gần nhau, như ghi nhớ kiến thức, giao tiếp, động lực và chi phí. Thông tin liên hệ của người trả lời không được đưa vào báo cáo.

Từ kết quả phân tích, nhóm xác định yêu cầu và thiết kế các hoạt động học tương ứng. Chẳng hạn, nhu cầu luyện giao tiếp được cụ thể hóa thành chức năng hội thoại theo tình huống; nhu cầu ghi nhớ và tìm lại kiến thức được thể hiện qua tra cứu, flashcard và dữ liệu học tập. Việc phát triển được chia thành từng phần để có thể kiểm tra sớm, nhận diện vướng mắc và điều chỉnh trước khi bổ sung chức năng tiếp theo.

Sau khi xây dựng, nhóm kiểm thử các hoạt động chính như đăng nhập, chọn bài, thực hiện bài luyện và xem lại kết quả. Với chức năng AI, việc đánh giá chú ý đến sự phù hợp của câu trả lời, mức độ dễ hiểu và thời gian chờ. Phản hồi dùng thử giúp xác định những bước còn khó sử dụng và nội dung cần cải thiện. Các kết quả được trình bày cùng điều kiện kiểm tra trong Chương 3 để làm rõ mức độ đáp ứng của sản phẩm.

\section{SCIENTIFIC AND PRACTICAL SIGNIFICANCE}

Ý nghĩa thực tiễn của NihonGO bắt đầu từ những hoạt động gần với nhu cầu tự học: có nơi tra lại từ vừa gặp, có bài để luyện, có cơ hội thử nói và có phản hồi để xem lại cách diễn đạt. Khi được tổ chức thuận tiện trong một hệ thống, những hoạt động này có thể giúp người học chủ động hơn trong việc ôn kiến thức và bố trí thời gian luyện tập. Đối với sinh viên, đây là hướng hỗ trợ bổ sung cho việc học trên lớp và những thời điểm chưa có người cùng thực hành.

Đề tài cũng tạo cơ hội vận dụng kiến thức về phát triển phần mềm và AI vào một bài toán cụ thể. Thay vì chỉ dừng ở việc kết nối một dịch vụ AI, nhóm cần giải quyết cách chuẩn bị nội dung, lưu dữ liệu, trình bày phản hồi và tổ chức quá trình học. Những kinh nghiệm từ quá trình xây dựng, kiểm thử và dùng thử có thể tiếp tục được sử dụng để cải thiện sản phẩm hoặc nghiên cứu cách hỗ trợ học ngoại ngữ bằng công nghệ.

Về định hướng khởi nghiệp, NihonGO có thể bắt đầu từ nhóm sinh viên có nhu cầu tự học, sau đó xem xét mở rộng học liệu và hợp tác với câu lạc bộ hoặc đơn vị đào tạo. Một phương án kinh doanh cần nghiên cứu là duy trì các chức năng cơ bản miễn phí và cung cấp gói nâng cao cho những hoạt động có chi phí AI lớn hơn. Trước khi lựa chọn mức giá và mở rộng, nhóm cần tiếp tục tìm hiểu người dùng có quay lại học thường xuyên hay không, chức năng nào được sử dụng nhiều và giá trị nào khiến họ sẵn sàng trả phí. Đây là các bước cần thiết để phát triển từ một đồ án kỹ thuật thành sản phẩm có khả năng duy trì lâu dài.

\section{THESIS STRUCTURE}

Báo cáo gồm phần Introduction và bốn chương. Phần Introduction giới thiệu lý do chọn đề tài, nhu cầu người học, mục tiêu và phạm vi thực hiện. Chương 1 trình bày cơ sở lý thuyết, tìm hiểu các sản phẩm liên quan và xác định yêu cầu hệ thống. Chương 2 mô tả kiến trúc, cơ sở dữ liệu, giao diện và cách xây dựng các chức năng của NihonGO. Chương 3 trình bày kết quả triển khai, kiểm thử, đánh giá và đề xuất mô hình kinh doanh. Chương 4 tổng hợp kết quả đạt được, các hạn chế và hướng phát triển tiếp theo.
